\chapter*{Glossary}
\addcontentsline{toc}{chapter}{Glossary}
\begingroup
\setlength{\LTleft}{0.04936\textwidth}
\setlength{\LTright}{0pt}
\begin{longtable}{@{}p{0.41159\textwidth}@{}p{0.53905\textwidth}@{}}
accelerator lattice & Ordered sequence of magnets, radiofrequency systems, drifts, and other elements that determines charged-particle transport through an accelerator.\\
attention & Neural-network mechanism that forms weighted combinations of value representations using learned similarities between queries and keys.\\
autoencoder & Neural network trained to encode data into a lower-dimensional representation and decode that representation back to the input space.\\
\noalign{\vskip\baselineskip}
beam centroid & First statistical moment of a beam-coordinate distribution; for coordinate \(\xi_{i} ,\) it is the ensemble average \(\langle \xi_{i} \rangle .\)\\
beam coupling impedance & Frequency-domain representation of the electromagnetic response of the accelerator environment to the beam.\\
beam size & Transverse spatial extent of a beam, conventionally expressed through an RMS coordinate spread and related to emittance, optical beta function, dispersion, and momentum spread.\\
betatron motion & Oscillatory transverse motion of a particle around the reference orbit under the action of focusing elements.\\
bunch & Localized ensemble of charged particles transported together through an accelerator.\\
bunch charge & Total electric charge carried by one particle bunch.\\
\noalign{\vskip\baselineskip}
causal attention & Attention mask that permits a token to receive information only from itself and preceding tokens in an ordered sequence.\\
collective effects & Beam-dynamics phenomena in which particles interact through fields generated by the beam itself or by its electromagnetic interaction with the accelerator environment.\\
conditional latent uncertainty & Variability obtained by repeatedly sampling the conditional latent distribution while holding model parameters and conditioning variables fixed.\\
conditional variational autoencoder & Variational autoencoder whose posterior and generation process are conditioned on auxiliary physical or machine variables.\\
convolution & Integral operation combining a function with a translated kernel; in beam dynamics it relates a longitudinal density to its induced wake potential.\\
coordinate-wise RMS spread & Square root of the second central moment of one beam coordinate; it remains defined for non-Gaussian distributions with finite variance but does not fully characterize their shape.\\
coupled-bunch instability & Coherent instability in which wake-mediated interactions couple the motion of different bunches circulating in an accelerator.\\
\noalign{\vskip\baselineskip}
\texttt{cpymad} & Python interface to MAD-X used to load and manipulate accelerator lattice descriptions.\\
\noalign{\vskip\baselineskip}
decoder & Model component that maps a latent representation back to the data or observable space.\\
digital twin & Computational representation of a physical system that combines models and data to support prediction, monitoring, calibration, or control.\\
discretization invariance & Property of an operator-learning architecture whose parameters do not depend on one fixed input or output discretization and whose discrete realizations converge under refinement.\\
dispersion & Dependence of the closed orbit or particle position on relative momentum deviation; it contributes to the measured transverse beam size.\\
\noalign{\vskip\baselineskip}
emittance & Phase-space area or volume characterizing the spread of a beam in a conjugate coordinate– momentum plane.\\
empirical measure & Discrete probability measure formed by assigning equal point mass to each simulated macroparticle.\\
encoder & Model component that maps high-dimensional observations into features or a latent representation.\\
\noalign{\vskip\baselineskip}
Fourier Neural Operator & Neural-operator architecture that parameterizes non-local integral transformations through learned multiplication of Fourier modes.\\
Fourier transform & Transformation between a function in physical space and its representation in frequency or wave-number space.\\
\noalign{\vskip\baselineskip}
graph message passing & Update mechanism in which graph nodes aggregate learned messages from neighbouring nodes connected by edges.\\
Graph Neural Operator & Neural-operator architecture that discretizes a learned kernel through graph-based message passing on arbitrary or physically structured nodes.\\
\noalign{\vskip\baselineskip}
Haïssinski distribution & Stationary longitudinal bunch distribution obtained from the self-consistent balance of RF confinement, wakefields, synchrotron-radiation damping, and quantum excitation below the microwave-instability threshold.\\
\noalign{\vskip\baselineskip}
inductive bias & Architectural assumption that favours particular solution structures, such as locality, causality, translation structure, or graph connectivity.\\
intrabeam scattering & Multiple small-angle Coulomb scattering among particles within the same bunch, which can increase beam emittances and energy spread.\\
\noalign{\vskip\baselineskip}
kernel & Function defining weighted interactions or an integral operator between input and output locations.\\
Kullback–Leibler divergence & Asymmetric divergence measuring how one probability distribution differs from another; in a VAE it regularizes the approximate posterior toward the prior.\\
\noalign{\vskip\baselineskip}
latent space & Learned lower-dimensional representation in which complex observations are encoded and manipulated.\\
Liouville’s theorem & Statement that Hamiltonian phase-space flow preserves phase-space volume and transports densities incompressibly.\\
local-window attention & Attention mechanism restricted to a finite neighbourhood of each token, imposing local support on information transfer.\\
longitudinal line density & Particle or charge density projected onto the longitudinal coordinate.\\
\noalign{\vskip\baselineskip}
macroparticle & Numerical simulation particle representing many physical particles while carrying effective charge and phase-space coordinates.\\
microwave instability & Single-bunch longitudinal instability associated with short-range wakefields and highfrequency impedance.\\
Monte Carlo estimator & Approximation of an expectation or integral by an average over random samples.\\
\noalign{\vskip\baselineskip}
neural operator & Learned parametric mapping between function spaces rather than only between fixeddimensional vectors.\\
\noalign{\vskip\baselineskip}
operator learning & Machine-learning framework for approximating mappings that transform input functions, fields, or distributions into output functions, fields, or distributions.\\
\noalign{\vskip\baselineskip}
Panofsky–Wenzel theorem & Relation connecting longitudinal and transverse wakefields through derivatives of the integrated electromagnetic force.\\
phase space & Space of particle positions and conjugate momenta; the beam state in this thesis is represented in six dimensions.\\
phase-space marginal & Lower-dimensional distribution obtained by integrating the full phase-space density over unobserved coordinates.\\
posterior distribution & Probability distribution updated after conditioning on observed data or model inputs.\\
prior distribution & Probability distribution specified before conditioning on observations; the standard VAE prior is commonly a unit Gaussian.\\
pushforward measure & Distribution obtained by transporting a measure through a deterministic map.\\
PyColleff & Python module developed by LNLS for impedance modelling and collectiveinstability analysis in electron storage rings.\\
PyWIT & Python Wake and Impedance Toolbox used to represent and manipulate accelerator wake and impedance models.\\
\noalign{\vskip\baselineskip}
quantile & Value below which a specified fraction of a probability distribution lies; useful for describing non-Gaussian beam tails and asymmetric distributions.\\
\noalign{\vskip\baselineskip}
reparameterization trick & Differentiable sampling construction used in variational autoencoders, typically written as a deterministic transformation of Gaussian noise.\\
resistive-wall impedance & Beam coupling impedance produced by finite electrical conductivity of the vacuum chamber walls.\\
\noalign{\vskip\baselineskip}
self-consistent force & Force that depends on the instantaneous beam distribution and therefore couples the evolution of individual particles to the collective state.\\
sliced Wasserstein distance & Distributional distance obtained by averaging one-dimensional Wasserstein distances over projected directions.\\
space charge & Collective electromagnetic interaction arising directly from the charge and current distribution of the beam.\\
surrogate model & Computationally inexpensive approximation of a more costly simulator, experiment, or physical model.\\
symplectic map & Phase-space transformation that preserves the symplectic structure of Hamiltonian dynamics and therefore respects the associated geometric invariants.\\
synchrotron radiation & Electromagnetic radiation emitted by relativistic charged particles undergoing transverse acceleration, producing damping and quantum excitation in electron rings.\\
\noalign{\vskip\baselineskip}
Transformer & Sequence-model architecture based on attention layers and token-wise transformations.\\
Twiss parameters & Optical functions \(\alpha , \beta ,\) and \(\gamma\) used to describe linear transverse beam envelopes and phase-space ellipses.\\
\noalign{\vskip\baselineskip}
uncertainty calibration & Agreement between predicted uncertainty intervals or probabilities and their empirical coverage or observed error frequencies.\\
under-dispersion & Condition in which the variability predicted by an ensemble or probabilistic model is systematically smaller than the observed prediction error.\\
\noalign{\vskip\baselineskip}
variational autoencoder & Probabilistic autoencoder that learns an approximate latent posterior and is trained using reconstruction and prior-regularization terms.\\
Vlasov equation & Kinetic equation describing collisionless evolution of a phase-space density under external and self-consistent forces.\\
Vlasov–Fokker–Planck equation & Kinetic equation combining Vlasov transport with damping and diffusion terms.\\
\noalign{\vskip\baselineskip}
wake function & Green-function response of the accelerator environment to a point-like source charge.\\
wake potential & Integrated electromagnetic effect experienced by a beam distribution, obtained by convolving the wake function with the source distribution.\\
Wasserstein distance & Optimal-transport distance between probability distributions that accounts for the cost of moving probability mass.\\
\noalign{\vskip\baselineskip}
Xsuite & Open-source Python suite for multiparticle accelerator simulations, including tracking, optics manipulation, collective-effects interfaces, and accelerator-model workflows.\\
Xtrack & Xsuite component used to construct accelerator lines and track macroparticles through lattice elements.\\
\end{longtable}
\endgroup
